% Drawing for measurement `mos-body-diode`
% Stands for the deck lines:
%   BS 0 s I=exp(2.302585092994046*V(x))
%   MD 0 g s s DUT / VG g s 0
%   .dc VX -6 2 0.05
%   L1 vr s {l} / SW s 0 ctrl 0 SWIDEAL / VR vr 0 30
\documentclass[border=8pt]{standalone}
\usepackage[siunitx,RPvoltages,european]{circuitikz}
\ctikzset{bipoles/length=0.9cm, label/align=straight, inductors/scale=0.8}
\ctikzset{inductor=american}
\usepackage{amsmath}
\begin{document}
\begin{circuitikz}[font=\footnotesize, line width=0.5pt]
\draw (0,0) node[nmos](Q){};
\node[anchor=west] at ($(Q.D)!0.5!(Q.S) + (0.55,0)$) {DUT};
\draw (Q.D) to[short] ++(0,0.3) to[short] ++(0,0.3) node[ground, rotate=180]{};
\draw (Q.S) to[short] ++(0,-0.5) coordinate(qs) to[isource, l_={\shortstack[r]{$I_S=10^{V(x)}$\\ \SI{1}{\micro A}\dots\SI{100}{A}}}] ++(0,-1.9) node[ground]{};
\draw (qs) to[short] ++(1.8,0) to[rmeter, t=V, l={$V_{SD}$}] ++(0,-1.9) node[ground]{};
\draw (Q.G) to[short] ++(-2.2,0) coordinate(gl) to[vsource, l_={\shortstack[r]{$V_{GS}$ 0, 4.5, \SI{10}{V}\\ (to the source)}}, invert] (gl |- qs) to[short] (qs);
\node[anchor=west, align=left] at (-4.6,3.3) {forward: the source current forced into the body diode, drain grounded};
\draw (9,0) node[nmos](R){};
\node[anchor=west] at ($(R.D)!0.5!(R.S) + (0.55,0)$) {DUT};
\draw (R.G) to[short] ++(-0.6,0) |- ($(R.S)+(0,-0.3)$);
\draw (R.D) to[short] ++(0,0.3) to[short] ++(0,0.3) node[ground, rotate=180]{};
\draw (R.S) to[short] ++(0,-0.3) coordinate(rs) to[L, l_={$L_{di/dt}$}] ++(0,-1.9) to[vsource, l_={$V_R$ \SI{30}{V}}, invert] ++(0,-1.4) node[ground]{};
\draw (rs) to[short] ++(1.6,0) to[isource, l={$I_F$ 1, \SI{10}{A}}] ++(0,-1.9) node[ground]{};
\draw (rs) to[short] ++(3.2,0) to[cute open switch, l={ideal switch}] ++(0,-1.9) node[ground]{};
\node[anchor=west, align=left] at (6.5,3.3) {recovery: $I_F$ through the diode, commutated at $di/dt$ = 100, 200, \SI{400}{A/\micro s}};
\end{circuitikz}
\end{document}
